\lysh{36570}{2}

\begin{alist}
    \item Γνωρίζουμε πως
    \[ \bar{x} = \frac{x_1 + x_2 + \dots + x_v}{\nu} = \frac{2 \cdot 14 + 18 + 20 + 2 \cdot 12}{6} \iff \bar{x} = \frac{90}{6} \iff \bar{x} = 15. \]
    \item Για να βρούμε τη διάμεσο του δείγματος τοποθετούμε όλες τις τιμές στη σειρά κατά αύξουσα σειρά. Δηλαδή, 12 12 14 14 18 20.

    Επειδή το δείγμα μας έχει άρτιο πλήθος παρατηρήσεων τη διάμεσο $ \delta $ την υπολογίζουμε με το ημιάθροισμα των μεσαίων βαθμών, δηλαδή, $ \delta = \frac{14+14}{2} = \frac{28}{2} = 14 $.

    Άρα, $ \delta=14 $.

    Το εύρος $R= \text{Μεγαλύτερος βαθμός} - \text{Μικρότερος βαθμός} = 20-12=8$.

    Άρα το εύρος του δείγματος είναι 8.
    \item Ο συντελεστής μεταβολής δίνεται από τον τύπο $ CV = \frac{s}{\bar{x}} $.

    Μετά από την απλή αντικατάσταση των τιμών που γνωρίζουμε έχουμε:
    \[ CV = \frac{s}{\bar{x}} \iff CV = \frac{3}{15} \iff CV = 0,20 \iff CV = 20\%. \]
\end{alist}
\vspace{6mm}
\noindent\rule{\textwidth}{0.4pt}
\vspace{6mm}

